\documentclass[10pt,a4paper]{amsart}
\usepackage{amsmath,amssymb,amsthm,graphicx}
% Retained minimal preamble for the Stacks Project chapter more-morphisms.tex (More on Morphisms).
% Verbatim excerpt of the chapter preamble: Xy-pic setup (native preamble lines 14-18) and the native
% theorem declarations with the mathematical macros (native preamble lines 154-217). Copied mechanically;
% nothing added, removed or redefined. The wrapper class and the amsmath/amssymb/amsthm/graphicx base
% packages are supplied by the assembler, so the class/externaldocument part of the native preamble is omitted.
% For commutative diagrams we use Xy-pic
\usepackage[all]{xy}

% We use 2cell for 2-commutative diagrams.
\xyoption{2cell}
% Theorem environments.
%
\theoremstyle{plain}
\newtheorem{theorem}[subsection]{Theorem}
\newtheorem{proposition}[subsection]{Proposition}
\newtheorem{lemma}[subsection]{Lemma}

\theoremstyle{definition}
\newtheorem{definition}[subsection]{Definition}
\newtheorem{example}[subsection]{Example}
\newtheorem{exercise}[subsection]{Exercise}
\newtheorem{situation}[subsection]{Situation}

\theoremstyle{remark}
\newtheorem{remark}[subsection]{Remark}
\newtheorem{remarks}[subsection]{Remarks}

\numberwithin{equation}{subsection}

% Macros
%
\def\lim{\mathop{\mathrm{lim}}\nolimits}
\def\colim{\mathop{\mathrm{colim}}\nolimits}
\def\Spec{\mathop{\mathrm{Spec}}}
\def\Hom{\mathop{\mathrm{Hom}}\nolimits}
\def\Ext{\mathop{\mathrm{Ext}}\nolimits}
\def\SheafHom{\mathop{\mathcal{H}\!\mathit{om}}\nolimits}
\def\SheafExt{\mathop{\mathcal{E}\!\mathit{xt}}\nolimits}
\def\Sch{\mathit{Sch}}
\def\Mor{\mathop{\mathrm{Mor}}\nolimits}
\def\Ob{\mathop{\mathrm{Ob}}\nolimits}
\def\Sh{\mathop{\mathit{Sh}}\nolimits}
\def\NL{\mathop{N\!L}\nolimits}
\def\CH{\mathop{\mathrm{CH}}\nolimits}
\def\proetale{{pro\text{-}\acute{e}tale}}
\def\etale{{\acute{e}tale}}
\def\QCoh{\mathit{QCoh}}
\def\Ker{\mathop{\mathrm{Ker}}}
\def\Im{\mathop{\mathrm{Im}}}
\def\Coker{\mathop{\mathrm{Coker}}}
\def\Coim{\mathop{\mathrm{Coim}}}

% Boxtimes
%
\DeclareMathSymbol{\boxtimes}{\mathbin}{AMSa}{"02}

%
% Macros for moduli stacks/spaces
%
\def\QCohstack{\mathcal{QC}\!\mathit{oh}}
\def\Cohstack{\mathcal{C}\!\mathit{oh}}
\def\Spacesstack{\mathcal{S}\!\mathit{paces}}
\def\Quotfunctor{\mathrm{Quot}}
\def\Hilbfunctor{\mathrm{Hilb}}
\def\Curvesstack{\mathcal{C}\!\mathit{urves}}
\def\Polarizedstack{\mathcal{P}\!\mathit{olarized}}
\def\Complexesstack{\mathcal{C}\!\mathit{omplexes}}
% \Pic is the operator that assigns to X its picard group, usage \Pic(X)
% \Picardstack_{X/B} denotes the Picard stack of X over B
% \Picardfunctor_{X/B} denotes the Picard functor of X over B
\def\Pic{\mathop{\mathrm{Pic}}\nolimits}
\def\Picardstack{\mathcal{P}\!\mathit{ic}}
\def\Picardfunctor{\mathrm{Pic}}
\def\Deformationcategory{\mathcal{D}\!\mathit{ef}}
\newenvironment{reference}{}{}
\newenvironment{slogan}{}{}
\newenvironment{history}{}{}
\usepackage{cleveref}
\crefname{originallemma}{original lemma}{original lemmas}
\Crefname{originallemma}{Original lemma}{Original lemmas}
\makeatletter
\newlabel{descent-lemma-locally-finite-nr-irred-local-fppf}{{descent-lemma-locally-finite-nr-irred-local-fppf}{source}{}{sourcepointer-1}{}}
\newlabel{descent-lemma-locally-finite-nr-irred-local-fppf@cref}{{[originallemma][0][]descent-lemma-locally-finite-nr-irred-local-fppf}{[1][0][]source}{}{}{}}
\makeatother
\crefname{originallemma}{original lemma}{original lemmas}
\Crefname{originallemma}{Original lemma}{Original lemmas}
\makeatletter
\newlabel{lemma-unramified-dominant-unibranch-is-etale}{{lemma-unramified-dominant-unibranch-is-etale}{source}{}{sourcepointer-2}{}}
\newlabel{lemma-unramified-dominant-unibranch-is-etale@cref}{{[originallemma][0][]lemma-unramified-dominant-unibranch-is-etale}{[1][0][]source}{}{}{}}
\makeatother
\crefname{originallemma}{original lemma}{original lemmas}
\Crefname{originallemma}{Original lemma}{Original lemmas}
\makeatletter
\newlabel{morphisms-lemma-dominant-finite-number-irreducible-components}{{morphisms-lemma-dominant-finite-number-irreducible-components}{source}{}{sourcepointer-3}{}}
\newlabel{morphisms-lemma-dominant-finite-number-irreducible-components@cref}{{[originallemma][0][]morphisms-lemma-dominant-finite-number-irreducible-components}{[1][0][]source}{}{}{}}
\makeatother
\title{A connected unramified cover of a geometrically unibranch integral scheme with a dominant component is etale and irreducible}
\author{The Stacks Project Authors}
\date{stacks-project:more-morphisms.tex:local-snapshot-20261011; GNU FDL 1.2 or later}
\begin{document}
\maketitle
\noindent\textit{Editorial note.} Statement span is lines 10495-10508 of sources/stacks-more-morphisms/pinned-original-source.tex; proof span is lines 10510-10527, the complete author proof (16 body lines).\par\medskip
\noindent\textit{Editorial note.} Navigation stubs: 3 cited-but-not-included labels are supplied as non-load-bearing pointers with their pinned tags where the snapshot tag index has them.\par\medskip
\noindent\textit{Editorial note.} extra\_packages is empty on purpose: the retained preamble already loads xy with options and a second declaration clashes in the assembler wrapper.\par\medskip
\noindent\textit{Editorial note.} Text is retained verbatim, including the author's own typos and placeholders.\par\medskip
\noindent\textit{Editorial note.} The author's statement contains a reference environment citing SGA 1, so the card carries a bibliography fragment. That fragment reproduces, field for field and without editing, the @BOOK\{SGA1\} entry from the Stacks Project snapshot bibliography shipped with this corpus at focused-seeds/AGP/sources/stacks-snapshot-20261011/my.bib (same upstream Stacks Project snapshot). No entry text was altered and no citation was reconstructed from memory.\par\medskip
\noindent\textit{Allowed standard prerequisites (editorial).}\begin{itemize}
\item Explicit prerequisite boundary. Unramified morphisms, geometrically unibranch schemes, and the flatness criterion for etaleness; the argument is retained.
\end{itemize}\medskip
\ifdefined\hypertarget\hypertarget{sourcepointer-1}{}\fi
\noindent\textit{Source reference.} Original lemma descent-lemma-locally-finite-nr-irred-local-fppf: sources/stacks-more-morphisms/pinned-original-source.tex, source lines 10527--10527. Cited by the retained statement or proof but not part of the collected spans; supplied only as a navigation pointer. The argument does not depend on its content.\par\medskip
\ifdefined\hypertarget\hypertarget{sourcepointer-2}{}\fi
\noindent\textit{Source reference.} Original lemma lemma-unramified-dominant-unibranch-is-etale: sources/stacks-more-morphisms/pinned-original-source.tex, source lines 10466--10475. Cited by the retained statement or proof but not part of the collected spans; supplied only as a navigation pointer. The argument does not depend on its content.\par\medskip
\ifdefined\hypertarget\hypertarget{sourcepointer-3}{}\fi
\noindent\textit{Source reference.} Original lemma morphisms-lemma-dominant-finite-number-irreducible-components: sources/stacks-more-morphisms/pinned-original-source.tex, source lines 10527--10527. Cited by the retained statement or proof but not part of the collected spans; supplied only as a navigation pointer. The argument does not depend on its content.\par\medskip
\begin{lemma}
\label{lemma-global-unramified-dominant-unibranch-is-etale}
\begin{reference}
\cite[Expose I, Corollary 9.11]{SGA1}
\end{reference}
Let $f : X \to Y$ be a morphism of schemes. Assume
\begin{enumerate}
\item $Y$ is integral and geometrically unibranch,
\item at least one irreducible component of $X$ dominates $Y$,
\item $f$ is unramified, and
\item $X$ is connected.
\end{enumerate}
Then $f$ is \'etale and $X$ is irreducible.
\end{lemma}

\begin{proof}
Let $X' \subset X$ be the irreducible component which dominates $Y$.
This means that the generic point of $X'$ maps to the generic point of $Y$
(see for example Morphisms, Lemma
\ref{morphisms-lemma-dominant-finite-number-irreducible-components}).
By Lemma \ref{lemma-unramified-dominant-unibranch-is-etale}
we see that $f$ is \'etale at every point of $X'$. In particular,
the open subscheme $U \subset X$ where $f$ is \'etale contains $X'$.
Note that every quasi-compact open of $U$ has finitely many irreducible
components, see
Descent, Lemma \ref{descent-lemma-locally-finite-nr-irred-local-fppf}.
On the other hand since $Y$ is geometrically unibranch and $U$ is \'etale
over $Y$, the scheme $U$ is geometrically unibranch. In particular, through
every point of $U$ there passes at most one irreducible component.
A simple topological argument now shows that $X' \subset U$
is both open and closed. Then of course $X'$ is open and closed in $X$ and by
connectedness we find $X = U = X'$ as desired.
\end{proof}

\begin{thebibliography}{1}
\bibitem{SGA1}
Grothendieck, Alexander.
\newblock {\em Rev\^{e}tements \'{e}tales et groupe fondamental (SGA 1)}.
\newblock Springer-Verlag, Lecture notes in mathematics 224, 1971.
\end{thebibliography}
\end{document}
